
quotient de @@M0@@ par le @@M1@@-ième facteur. Dans ce système de morphismes, les quotients canoniques se correspondent donc par image inverse. Le système inductif ainsi défini sera noté @@M2@@.

On définit son @@M3@@ par la formule :
@@M4@@

Soient @@M5@@ les faisceaux canoniques sur @@M6@@ (4.5), et @@M7@@ leur classe dans @@M8@@. Posons @@M9@@. Si on fait varier @@M10@@, les @@M11@@ se correspondent par les morphismes de transition, et définissent par conséquent des éléments, encore notés @@M12@@, de @@M13@@.

\textbf{Proposition 4.10.} \uline{Soit} @@M14@@ \uline{un schéma quasi}-\uline{compact}. \uline{Alors}, \uline{avec les notations de} 4.9, \uline{on a un isomorphisme}~:
@@M15@@

D'après 4.8, @@M16@@ est la @@M17@@-algèbre engendrée par les @@M18@@, @@M19@@, et les @@M20@@, @@M21@@, soumis aux relations déduites de :
@@M22@@
soit encore :
@@M23@@

Cette relation permet d'exprimer les @@M24@@ en fonction des @@M25@@, et donne une série de relations isobares entre les @@M26@@ (@@M27@@ étant supposé de poids @@M28@@) en écrivant que les coefficients de @@M29@@ sont nuls en degré @@M30@@; de plus, ces relations engendrent l'idéal des relations entre les @@M31@@. On peut donc écrire :
